% -*- coding: utf-8 -*-
% !TeX program = xelatex
% !BIB program = biber

\documentclass{nktba}
\usepackage{nktbacs}

\usepackage{booktabs} % 三线表
\usepackage{graphicx} % 图片导入
\usepackage{hologo} % logo
\usepackage{xurl}

% 文档元数据
\title{南开大学本科生毕业论文模板 \\ 计网版}
\title*{version: 2023/06/02}
\studentid{1234567}
\author{XXX}
\grade{2019 级}
\major{计算机科学与技术}
\department{计算机科学与技术}
\college{计算机学院}
\adviser{XXX\quad 教授 \\ 罗翔\quad 教授 \\（张三大学）}
\date{\today}
\keywords{南开大学；毕业论文；LaTeX模板}
\keywords*{dark energy model; dynamic analysis}
\addbibresource{reference.bib} % 参考文献

\begin{document}

\maketitle

\begin{abstract}
摘要内容 TODO
\end{abstract}
  
\begin{abstract*}
Since 1998, two independent supernova research groups have discovered that the universe is accelerating, and the dark energy has become a hot topic in cosmology. 
\end{abstract*}

\tableofcontents

\section{介绍}

\subsection{注意事项}

提供本模板的目的是为了给您的论文写作带来方便，然而，
作者不保证这个模板完全符合学校的要求，也不对由此产生的任何后果负责。请使用者认真校对模板格式是否符合各学院的格式要求，确保及时发现问题，以便模板维护者尽早解决。

\section{关于模板}

\subsection{求助与反馈}

请follow我们最新的进度\footnote{\url{https://github.com/Tr0py/NKU-thesis-template-2020}}, 
欢迎加入毕业论文模板维护的工作中。如果遇到使用问题，你可以用以下方式求助：
\begin{itemize}
  \item 在 \href{https://github.com/Tr0py/NKU-thesis-template-2020/discussions}{GitHub讨论区} 询问你不懂的地方。
  \item 在 \href{https://github.com/Tr0py/NKU-thesis-template-2020/issues}{GitHub Issue} 中提出你遇到的待修复的问题。
  \item 如果不会使用GitHub，你可以在QQ群中提问：群号469868290。
\end{itemize}

\subsection{贡献者}

感谢历届的贡献者们：
\begin{itemize}
  \item 赵子懿，20届本科毕业生，\url{troppingz@gmail.com}，\href{https://tr0py.github.io/}{主页}
  \item 潘宇，20届本科毕业生，\url{panyu@nbjl.nankai.edu.cn}，\href{https://nbjl.nankai.edu.cn/2019/0513/c19244a264683/page.htm}{主页}
  \item 黎鸿儒，22届本科毕业生，\url{lihongru@mail.nankai.edu.cn}
  \item 林雪，22届本科毕业生，\url{ylsnowy@gmail.com}
  \item 费迪，22届本科毕业生，\url{fd@nbjl.nankai.edu.cn}
  \item 安祺，23届本科毕业生，\url{aqnin@outlook.com}
\end{itemize}

\subsection{发展历史}

\subsubsection{发源}

\LaTeX 由于其专业的排版效果，擅长处理复杂的数学公式，免费等优势而得到了广泛的使用，成为许多西文期刊接收投稿的主流格式；而\LaTeX 环境的安装，编辑器的选择以及编译\LaTeX 文件的设置等增加了使用者的使用困难，并且在此过程中出现的问题有难有易，种类繁多，不易找到解决方案。Overleaf作为一档在线\LaTeX 编辑网站，使用简单，支持在线和跨平台编译修改毕业论文，并且支持多人协作以方便老师和项目组其他人修改。因此，我们参考南开大学08级本科生卓小杨依据《南开大学本科生毕业论文（设计）管理规定（2011年 修订）》修改制成的本科生毕业论文模板，将南开大学孙文昌老师制作的南开大学研究生论文\LaTeX 模板移植到了Overleaf上，借鉴未知人士对其的修改，针对该模板的若干问题和版本旧的不足进行了修改及更新。

本模板参考《南开大学本科生学位论文排版规范（试行）(2016 修正）版本 1.2.21》修改而成，更新了封面格式、添加了所需字体、更新参考文献到最新的bibtex标准。作者希望能给使用者写作论文带来方便。然而，作者不保证本模板完全符合学校要求，也不对由此带来的风险和损失承担任何责任。如果您不接受这项条款，请不要使用本模板。

\hfill 16级本科生\quad 赵子懿\quad 潘宇

\subsubsection{迭代}

在2023年5月，本模板经历了一次重构。根据 \href{http://jwc.nankai.edu.cn/2022/1124/c24a497818/page.htm}{2022年11月版《南开大学本科毕业论文（设计）指导手册》} 的格式要求重新进行了实现，不再兼容原模板使用接口。

\hfill 19级本科生\quad 安祺

\section{使用说明} \label{chpt:A}

本模板基于 ctexart 实现，如本文介绍的内容不够详细丰富，可查看 ctex 宏基文档了解更多信息。

\subsection{入门资料}

如果你需要相对全面的了解 latex 的使用，推荐阅读\href{http://mirrors.ctan.org/info/lshort/chinese/lshort-zh-cn.pdf}{《一份（不太）简短的 LATEX 2ε 介绍》}。如果你想要学习具体问题的解决方法，推荐阅读 \href{https://www.overleaf.com/learn}{Overleaf文档}。

\subsection{字体}

本模板在不同平台默认使用的字体不尽相同，如\ref{tab:ctex-fontset}所示。如果需要修改 ctex 的 fontset 选项，可以在引入本模板文档类前向ctexart 传入参数，例如使用 \verb| \PassOptionsToClass{fontset=windows}{ctexart}| 显示使用windows字体配置。不过，你需要确保这些字体在你的系统上可用。

\begin{table}
  \centering
  \caption{不同平台的默认字体}
  \begin{tabular}{ccc}
  \toprule
  平台 & 中文字体 & 英文字体 \\
  \midrule
  Windows & 中易宋黑楷仿 & Times New Roman \\
  Overleaf & Fandol宋黑楷仿 & Times New Roman\\
  Other Linux & Fandol宋黑楷仿 & STIX \\
  \bottomrule
  \end{tabular}
  \label{tab:ctex-fontset}
\end{table}

有时，ctex宏包会使用黑体作为宋体的粗体使用，你可以使用本模板提供的\verb|\textfb| 来进行伪加粗。一般情况下您不需要显式地设置文本字体。如果确实需要，请考虑使用以下命令：

\begin{description}
\item[宋体] \verb|\songti| 或 \verb|\textnormal{<文本>}|
\item[黑体] \verb|\heiti| 或 \verb|\textsf{<文本>}|
\item[楷体] \verb|\kaishu|  或 \verb|\textit{<文本>}|
\item[仿宋] \verb|\fangsong| 或 \verb|\texttt{<文本>}|
\item[粗体] \verb|\bfseries| 或 \verb|\textbf{<文本>}|
\item[伪加粗] \verb|\textfb[<加粗程度>]{<文本>}|
\end{description}

使用 \hologo{XeLaTeX}\ 编译时，会自动在中英文转换时添加必要的空格。但是某些情况中，空格并不能被自动添加。比如，\verb|中文\url{https://google.com}| 的显示效果为：中文\url{https://google.com}。这时就需要手动添加，这类情况的特点是内容被一些命令所包裹。

\subsection{交叉引用} \label{sec:ex:A}

修改 \verb|\the<counter>| 的定义，可以修改编号的显示格式。 chngcntr 宏包提供了 \verb|\counterwithin{<counter>}{<withincounter>}| 宏对此进行了封装。 该命令会将 \verb|<counter>| 关联到 \verb|<withincounter>|，并将 \verb|\the<counter>| 的定义修改为 \verb|\the<within>.\arabic{<counter>}|。可以通过该命令实现按章编号，如\ref{lst:counterwithin}所示。

\begin{lstlisting}[
	caption={counterwithin 命令使用示例},
	label={lst:counterwithin},
	language={[LaTeX]TeX},
	morekeywords={counterwithin}
]
\counterwithin{table}{section}
\counterwithin{figure}{section}
\counterwithin{equation}{section}
\end{lstlisting}

\LaTeX 内核内置了 \verb|\labelformat{<counter>}{<format>}| 宏来设置交叉引用的标签格式，使用示例如\ref{lst:labelformat}所示。hyperref提供的 \verb|\autoref| 也可以实现类似的功能，但是只能修改前缀。

\begin{lstlisting}[
	caption={labelformat 命令使用示例},
	label={lst:labelformat},
	language={[LaTeX]TeX},
	morekeywords={labelformat,zhnumber},
  emph={\#1}
]
\labelformat{section}{第\zhnumber{#1}章}
\labelformat{figure}{图~#1\space}
\labelformat{table}{表~#1\space}
\end{lstlisting}

引用章节号可参考如下格式: \ref{chpt:A}\ref{sec:ex:A}。

\subsection{参考文献} \label{manual:ref}

由于各专业所需的参考文献格式不尽相同，本模板未统一配置参考文献样式。请使用者自行配置所需的参考文献样式。

当前主要使用 natbib 或 biblatex 宏包设置参考文献和引用格式，示例目前使用 biblatex 和内置样式排版，下面给出一个使用 natbib 排版 gbt7714-2015 样式的例子（需要配合bibtex使用），如\ref{lst:natbib-gbt7714}所示。如果你需要对更多细节进行调整，请查阅 natbib 或 biblatex 宏包文档。

\begin{lstlisting}[
	caption={用 natbib 排版 gbt7714-2015 样式的参考文献},
	label={lst:natbib-gbt7714},
	language={[LaTeX]TeX},
	morekeywords={labelformat,zhnumber},
  emph={\#1}
]
% 导言区
\usepackage[sort&compress]{gbt7714}
\bibliographystyle{gbt7714-numerical} % gbt7714-2015
\usepackage{natbib}

% 参考文献排版处
\begin{reference}
  \renewcommand{\bibsection}{}
  \bibliography{reference.bib}
\end{reference}
\end{lstlisting}

参考文献测试样例:
\cite{ChenCheChen2001,Nadkarni-1992,Hua-Wang-1973}
文献\cite{ZhuKeZhen,Huo}文献\cite{JiangXiZhou,Timoshenko,Zhang-Wang}
\cite{ZhuKeZhen}
测试\citep{PQEMU}或者\citep{li1988ivy}
第一次编译\cite{hong2012hqemu}\cite{zaharia2010spark}\cite{wenwenasplos2018}

\subsubsection{录入参考文献}

本模板采用 biblatex 宏包管理参考文献。如果你对此不熟悉，可以参考biblatex使用手册。

\subsection{图表}

插图示例，如\ref{fig:nku}所示。

\begin{figure}
\centering
\includegraphics[viewport=0 0 2984 969,width=40mm]{figures/nankaidaxue.pdf}
\caption{南开大学}
\label{fig:nku}
\end{figure}

\subsection{列表}

本节示例来自 \url{https://www.overleaf.com/learn/latex/Lists}

\begin{itemize}
  \item This is my first point.This is my first point.This is my first point.This is my first point.This is my first point.
  \item Another point I want to make 
  \item[!] A point to exclaim something!
  \item[NOTE] This entry has no bullet
  \item[] A blank label?
\end{itemize}

\begin{enumerate}
  \item This is my first point.This is my first point.This is my first point.This is my first point.This is my first point.
  \item Another point I want to make 
  \item[!] A point to exclaim something!
  \item[NOTE] This entry has no bullet
  \item[] A blank label?
\end{enumerate}

\begin{description}
  \item This is an entry \textit{without} a label.
  \item[Something short] A short one-line description.
  \item[Something long] A much longer description.A much longer description.A much longer description.A much longer description.A much longer description.A much longer description.A much longer description.A much longer description.A much longer description.A much longer description.A much longer description.A much longer description.A much longer description.A much longer description.A much longer description.A much longer description.A much longer description.A much longer description.A much longer description.A much longer description.A much longer description.A much longer description.A much longer description.A much longer description.A much longer description.A much longer description.A much longer description.
\end{description}

\begin{enumerate}
  \item First level item
  \item First level item
  \begin{enumerate}
    \item Second level item
    \item Second level item
    \begin{enumerate}
      \item Third level item
      \item Third level item
      \begin{enumerate}
        \item Fourth level item
        \item Fourth level item
      \end{enumerate}
    \end{enumerate}
  \end{enumerate}
\end{enumerate}

\begin{itemize}
  \item First level item
  \item First level item
  \begin{itemize}
    \item Second level item
    \item Second level item
    \begin{itemize}
      \item Third level item
      \item Third level item
      \begin{itemize}
        \item Fourth level item
        \item Fourth level item
      \end{itemize}
    \end{itemize}
  \end{itemize}
\end{itemize}

\subsection{一些建议}
\subsubsection{关于分数的写法}

\LaTeX 提供宏命令\verb+\frac+, 用以打印分数. 为使得版面整齐, 该命令的使用应遵循以下原则:

\begin{enumerate}
\item 仅在分行表达式中使用,
\item 不嵌套使用,
\item 不在上下标中使用.
\end{enumerate}

也就是说, 行内表达式和上下标中出现分数时一律用 $a/b$表示, 如
$(x+2)/((3x^2+4)(7+y))$. 下面是居中表达式:

\[
 x^2 = y^{1/2} +3.
\]

多行表达式: 尽量在加、减、乘、等号前换行. 在乘号前换行时,
下一行首用 \string\times:
\def\iint{\mathop{\int\!\!\!\int}}\def\calG{\mathcal G}
\begin{eqnarray}
&&\left|(W_{\psi_1}f)(a,b)-(W_{\psi_1}f)(a_j,b_{j,k})\right|^{2}\nonumber\\
&=&\frac{1}{C^{2}_{\varphi}}\Bigg|\iint_{\calG} (W_{\varphi}f)(s,t) \nonumber\\
&&\qquad\times \Bigg( (W_{\psi_1}\varphi)\left(\frac{a}{s},
\frac{b-t}{s}\right)
     -(W_{\psi_1}\varphi)\left(\frac{a_{j}}{s}, \frac{b_{j,k}-t}{s}\right)\Bigg)
  \frac{dsdt}{s^{d+1}}\Bigg|^2 \nonumber\\
&\le& \frac{1}{C^2_{\varphi}} \iint_{\calG} |(W_{\varphi}f)(s,t)|^2 \nonumber\\
&&\qquad \times\left| (W_{\psi_1}\varphi)\left(\frac{a}{s},
\frac{b-t}{s}\right)
    -(W_{\psi_1}\varphi)\left(\frac{a_{j}}{s}, \frac{b_{j,k}-t}{s}\right)\right|
   \frac{dsdt}{s^{d+1}}  \nonumber\\
&&\qquad \times   \iint_{\calG}\!
 \left|(W_{\psi_1}\varphi)\left(\frac{a}{s}, \frac{b-t}{s}\right)
    -(W_{\psi_1}\varphi)\left(\frac{a_{j}}{s}, \frac{b_{j,k}-t}{s}\right)\right|
 \frac{ ds dt}{s^{d+1}} \nonumber\\
&=& \frac{1}{C^2_{\varphi}} ....  \label{eq:a0}
\end{eqnarray}


\subsubsection{标点}
科技文献中一般用半角标点, 请参考《中国科学》发表的论文.

如果使用全角标点, 可以使用
\begin{verbatim}
  \punctstyle{<style>}
\end{verbatim}
选择标点样式, 有效值为
\begin{verbatim}
  quanjiao (所有标点符号占一个汉字宽度，
            相邻两个标点占一个半汉字宽度)
  banjiao  (所有标点符号占半个汉字宽度)
  hangmobanjiao (所有标点符号占一个汉字宽度，行末行首半角)
  kaiming  (句号、叹号、问号占一个汉字宽度，其他标点占半个汉字宽度)
\end{verbatim}
缺省为全角式。注意：不论选择哪种样式，都提供行末对齐(margin kerning)功能。


引用:   \ref{thm:latex}的推论是什么呢？
方程式编号:  由\ref{eq:a1}和\ref{eq:aa2}.



\section{常用宏包使用示例}

\subsection{公式定理}

\begin{theorem} \label{thm:latex}
\LaTeX 的输出是最完美的.
\end{theorem}

先证明一个引理
\begin{lemma} \label{thm:tex}
\TeX 文件在不同操作系统下的排版结果完全一致.
\end{lemma}

单个带编号的表达式
\begin{equation}\label{eq:a1}
x=y+z
\end{equation}

单个不带编号的表达式
\[
y=x-z.
\]

不带编号的多行表达式
\begin{eqnarray*}
x&=&y+z \\
  &=&z-s\\
  &<& 3. \\
  && \mbox{一些注释}
\end{eqnarray*}

带编号的多行表达式
\begin{eqnarray}
  x&=& y-z, \label{eq:aa1}\\
  y&=& x+z, \nonumber \\
  z&=&y-x. \label{eq:aa2}
\end{eqnarray}

\subsection{伪代码}

TODO，推荐阅读：\url{https://www.overleaf.com/learn/latex/Algorithms}

\subsection{代码块}

TODO，推荐阅读：\url{https://www.overleaf.com/learn/latex/Code_listing}

\appendix
\section{第一个附录}
此处展示的为单标题附录，你可以使用 \verb|\appendices| 来排版多个附录。

\begin{reference}
\printbibliography[heading=none]
\end{reference}

\begin{acknowledgement}

\end{acknowledgement}

\end{document}
